\documentclass[12pt]{article}
\usepackage{amsmath, amssymb, geometry}
\geometry{margin=1in}
\setlength{\parindent}{0pt}

\title{Quiz 4\\ \vspace{8pt} \Large ECON 308 -- International Economic Relations}
\author{Instructor: Muhebullah Karimzada}
\date{October 28,  2025}
\usepackage{comment}

\begin{document}
\maketitle

\section*{Multiple Choice Questions (1 mark each)}

\textbf{1.} If a country’s terms of trade improve, its welfare will: \\[3pt]
A. Always increase. \\
B. Always decrease. \\
C. Increase if it is an exporter of the good whose relative price rises. \\
D. Decrease if it imports the good whose relative price rises. \\[10pt]

\textbf{2.} A country experiences growth biased toward its export good. This tends to: \\[3pt]
A. Improve its terms of trade. \\
B. Deteriorate its terms of trade. \\
C. Leave its terms of trade unchanged. \\
D. Shift its PPF inward. \\[10pt]

\textbf{3.} Suppose the world relative price of a country’s export good rises due to increased global demand. According to the Standard Trade Model, this will: \\[3pt]
A. Improve the country’s terms of trade and raise its welfare. \\
B. Deteriorate the terms of trade and lower welfare. \\
C. Leave both production and consumption unchanged. \\
D. Make the country worse off even though export prices rise. \\[10pt]

\textbf{4.} If the world relative price of cloth rises, a country that exports cloth will experience: \\[3pt]
A. A deterioration in its terms of trade. \\
B. An improvement in its terms of trade. \\
C. No change in welfare. \\
D. A fall in export earnings. \\[10pt]
\newpage
\textbf{5.} In the Standard Trade Model, changes in a country’s welfare can be illustrated by: \\[3pt]
A. Movements along the PPF only. \\
B. The slope of the indifference curve only. \\
C. Tangency between the highest possible indifference curve and the trade line. \\
D. The area between import and export demand curves. \\[14pt]



\textbf{Problem 1 (7 marks)} \\[3pt]
Country A produces \textbf{food (F)} and \textbf{manufactures (M)} with the PPF given by:
\[
F^2 + M^2 = 100
\]
The world relative price of manufactures is \( P_M / P_F = 2 \). \\
Country A consumes at the point where its indifference curve satisfies \( M = 0.5F \). \\[6pt]

\textbf{Questions:}
\begin{enumerate}
    \item Find the production point \((F, M)\) given the relative price.
    \item Compute the consumption point assuming the country exports the good it produces in excess of consumption.
    \item Calculate the country’s terms of trade and illustrate (conceptually) whether it gains from trade compared to autarky.
\end{enumerate}

\vspace{1em}

\textbf{Problem 2 (8 marks)} \\[3pt]
Country B’s PPF before growth is:
\[
Q_T = 200 - 2Q_C
\]
where \( Q_T \) is textiles and \( Q_C \) is chemicals. The world relative price is \( P_T / P_C = 1 \). \\[4pt]
After biased growth toward textiles, the PPF becomes:
\[
Q_T = 240 - 2Q_C
\]

\textbf{Questions:}
\begin{enumerate}
    \item Compute the change in maximum production possibilities for both goods.
    \item Assuming world prices are fixed initially, find the change in export potential.
    \item Explain, using the Standard Trade Model framework, how this biased growth could affect the \textbf{terms of trade} and \textbf{welfare} of Country B.
\end{enumerate}

\end{document}
